\section{Related Work}

This section groups previous work into five areas: Indian digital policing infrastructure, blockchain for digital evidence and audit, access control and privacy, explainable AI in high-stakes law-enforcement settings, and India-specific crime analytics. The purpose is to show what already exists and where the proposed work is positioned.

\subsection{Indian Digital Policing: CCTNS and ICJS}

CCTNS and ICJS provide the institutional baseline for this research. Official CCTNS information shows that India already has large-scale police digitization, including FIR and chargesheet digitization, data replication, search facilities, standardized investigation forms, and master codes \cite{pib_cctns_2026}. ICJS extends this by linking police/CCTNS with courts, prisons, forensics, and prosecution \cite{mha_icjs_2026}. These systems show that the Indian context is not a missing-infrastructure problem.

For this paper, the important observation is that official infrastructure descriptions do not provide an open research workload for access-request decisions, explanation traces, audit attacks, or policy evaluation. Therefore, the proposed work is not a replacement for CCTNS or ICJS. It is an overlay research model for access governance, auditability, and explanation.

\subsection{Blockchain for Evidence and Audit}

Blockchain has been widely studied as a tool for evidence integrity and chain-of-custody management. Kim et al. proposed a two-level blockchain system for digital crime evidence management, separating frequently changing investigation information from less frequently changing evidence such as videos \cite{kim_two_level_2021}. LEChain proposed a blockchain-based lawful evidence management scheme for digital forensics, including evidence access control and privacy concerns for witnesses and jurors \cite{li_lechain_2021}. A Hyperledger Fabric-based digital evidence management model also shows how permissioned blockchain can support distributed evidence management among authorized participants \cite{jeong_fabric_evidence_2020}.

These works support the idea that blockchain can improve tamper-evident audit and evidence provenance. However, they mostly focus on evidence lifecycle management, digital forensics, or chain-of-custody. The present work differs by focusing on inter-agency access governance: who should be allowed to access which sensitive record, under which policy, with what approval, and with what explanation.

Hyperledger Fabric is relevant because it is designed for permissioned networks with known organizations, membership identities, modular consensus, and configurable execution \cite{androulaki_fabric_2018}. NISTIR 8403 and blockchain access-control surveys show that blockchain can support auditability and tamper resistance, but also introduce governance, privacy, scalability, and policy-complexity issues \cite{nist_blockchain_access_2022,rouhani_access_control_survey_2019}. These sources prevent overclaiming: blockchain is useful as an audit layer, but it is not a complete privacy or authorization solution.

\subsection{Access Control, Privacy, and Policy Enforcement}

Access control is central to SEBA-XAI. NIST SP 800-162 defines ABAC as authorization based on subject, object, action, and environmental attributes evaluated against policies \cite{nist_abac_2014}. This is suitable for inter-agency police data sharing because access decisions depend on more than a user's broad role. A requester may have the right rank but not the correct jurisdiction, case assignment, approval token, or purpose.

Fabric plus ABAC has already been studied for fine-grained data sharing on Hyperledger Fabric \cite{zhao_fabric_abac_2022}. This prior work shows that blockchain and ABAC can be combined with encrypted off-chain storage and on-chain references. More recent accountable and privacy-preserving blockchain access-control work studies privacy and accountability concerns such as permission invisibility and access anonymity \cite{hu_accountable_private_ac_2024}. Privacy-preserving machine learning surveys also map techniques such as differential privacy, federated learning, secure multiparty computation, homomorphic encryption, and trusted execution environments \cite{xu_ppml_survey_2021,abadi_dp_sgd_2016}.

The limitation for the present research is that these works do not directly evaluate police-specific access-request workflows with superior approval, explanation artifacts, revocation cases, metadata exposure, and adversarial audit attacks. SEBA-XAI uses the access-control literature as a technical foundation, but the contribution is in integrating policy enforcement with audit and explanation in a controlled police-record-sharing benchmark.

\subsection{XAI, Fairness, and High-Stakes Law Enforcement}

Explainability is important in this domain because access decisions may affect privacy, investigations, and accountability. Rudin argues that high-stakes applications should prefer interpretable models where possible rather than relying only on explanations of black-box models \cite{rudin_2019}. Law-enforcement XAI literature also emphasizes that explanations must account for stakeholder needs, AI literacy, and automation-bias risk \cite{zocholl_xai_law_enforcement_2025}.

Criminal-justice AI literature warns against overclaiming predictive systems. Dressel and Farid showed limits of complex recidivism prediction tools compared with simpler baselines \cite{dressel_fairness_limits_2018}. Fairness research shows that different fairness criteria can conflict, especially when base rates differ across groups \cite{chouldechova_2017,kleinberg_tradeoffs_2016}. Ensign et al. showed that predictive-policing systems can create feedback loops when policing attention influences observed future data \cite{ensign_feedback_2018}. These works motivate a conservative research direction: this paper does not attempt individual suspect prediction or crime prediction. It focuses on access governance and explanation of access decisions.

Some recent works combine blockchain, XAI, and legal or justice contexts. Blockchain-based auditing of legal decisions supported by explainable AI and generative AI tools connects audit trails with AI-supported legal decision-making \cite{sachan_legal_ai_blockchain_2024}. A blockchain-based XAI-justice system also shows that privacy, blockchain, and XAI can be combined at an architecture level \cite{demertzis_xai_justice_2023}. These works are close in spirit, but they are not designed around CCTNS/ICJS-compatible inter-agency police record access.

\subsection{India-Specific Crime Analytics and Legal Data}

India-specific crime analytics commonly uses aggregate NCRB data. Studies have analyzed spatial and demographic crime patterns, including crime against women and murder-motive forecasting with XAI methods \cite{sharma_india_crime_2023,naick_murder_xai_2024,reza_crix_2025}. Indian legal-NLP datasets such as IndianBailJudgments-1200 are also emerging \cite{deshmukh_bail_2025}. These datasets and papers are useful for understanding Indian crime and legal-data research, but they do not solve the access-governance problem.

The key limitation is that public NCRB-style data is aggregate reported/registered crime data, not a record-level police access-control dataset \cite{ncrb_2023}. Therefore, SEBA-XAI uses synthetic access requests for the first benchmark. This keeps the work reproducible without making unsupported claims about restricted police records.

\subsection{Research Gap}

The reviewed work shows that the individual components already exist: CCTNS/ICJS infrastructure, blockchain evidence management, ABAC and PBAC-style access control, privacy-preserving mechanisms, and XAI for high-stakes decisions. The remaining gap is their integration and evaluation for a sensitive inter-agency police-record access workflow. Specifically, prior work does not sufficiently address a CCTNS/ICJS-compatible overlay that jointly evaluates contextual access decisions, superior-review escalation, off-chain record commitments, permissioned blockchain-style audit, explanation artifacts, metadata exposure, and adversarial audit cases. This is the gap addressed by SEBA-XAI.
